\hypertarget{hux1ea1-tux1ea7ng-dux1b0ux1edbi-dux1ea1ng-code-terraform}{%
\section{Hạ tầng dưới dạng Code
(Terraform)}}

Phần này mô tả việc định nghĩa và triển khai toàn bộ hạ tầng AWS cho
News RAG Pipeline bằng Terraform.

\hypertarget{kiux1ebfn-truxfac-terraform}{%
\subsection{Kiến trúc Terraform}}

File chính: \texttt{main.tf} --- định nghĩa tất cả tài nguyên AWS:

\begin{verbatim}
VPC & Network $\rightarrow$ Security Groups $\rightarrow$ RDS Aurora $\rightarrow$ ECR $\rightarrow$ ECS Cluster $\rightarrow$ IAM Roles $\rightarrow$ CloudWatch Logs $\rightarrow$ Task Definitions $\rightarrow$ EventBridge Scheduler
\end{verbatim}

\hypertarget{vpc-network-main.tf40-80}{%
\subsection{\texorpdfstring{1. VPC \& Network
(\texttt{main.tf:40-80})}{1. VPC \& Network (main.tf:40-80)}}}

\emph{(Mã nguồn chi tiết đã được lược bỏ để báo cáo ngắn gọn. Vui lòng
xem trong source code dự án)}

\textbf{Kết quả:} VPC \texttt{newsrag-vpc} với 2 public subnet ở 2 AZ,
Internet Gateway, Route Table cho phép traffic ra Internet.

\hypertarget{security-groups-main.tf82-114}{%
\subsection{\texorpdfstring{2. Security Groups
(\texttt{main.tf:82-114})}{2. Security Groups (main.tf:82-114)}}}

\emph{(Mã nguồn chi tiết đã được lược bỏ để báo cáo ngắn gọn. Vui lòng
xem trong source code dự án)}

\hypertarget{aurora-postgresql-serverless-v2-main.tf116-144}{%
\subsection{\texorpdfstring{3. Aurora PostgreSQL Serverless v2
(\texttt{main.tf:116-144})}{3. Aurora PostgreSQL Serverless v2 (main.tf:116-144)}}}

\emph{(Mã nguồn chi tiết đã được lược bỏ để báo cáo ngắn gọn. Vui lòng
xem trong source code dự án)}

\textbf{Lưu ý quan trọng:} - \texttt{db.t4g.medium} = 2 ACU (Aurora
Capacity Units) \textasciitilde{} \$15-20/tháng -
\texttt{skip\_final\_snapshot\ =\ true} chỉ dùng cho dev/test.
Production nên bật snapshot. - Sau khi deploy, chạy
\texttt{database/warehouse.sql} để tạo schema Star Schema + pgvector.

\hypertarget{ecr-repository-main.tf147}{%
\subsection{\texorpdfstring{4. ECR Repository
(\texttt{main.tf:147})}{4. ECR Repository (main.tf:147)}}}

\begin{Shaded}
\begin{Highlighting}[]
\NormalTok{resource "aws\_ecr\_repository" "api" \{ name = "newsrag{-}api" \}}
\end{Highlighting}
\end{Shaded}

Chứa Docker image cho cả Crawler, ETL, và Vectorize tasks.

\hypertarget{ecs-cluster-main.tf150}{%
\subsection{\texorpdfstring{5. ECS Cluster
(\texttt{main.tf:150})}{5. ECS Cluster (main.tf:150)}}}

\begin{Shaded}
\begin{Highlighting}[]
\NormalTok{resource "aws\_ecs\_cluster" "cluster" \{ name = "newsrag{-}cluster" \}}
\end{Highlighting}
\end{Shaded}

\hypertarget{iam-roles-main.tf152-180}{%
\subsection{\texorpdfstring{6. IAM Roles
(\texttt{main.tf:152-180})}{6. IAM Roles (main.tf:152-180)}}}

\emph{(Mã nguồn chi tiết đã được lược bỏ để báo cáo ngắn gọn. Vui lòng
xem trong source code dự án)}

\hypertarget{cloudwatch-log-group-main.tf183-186}{%
\subsection{\texorpdfstring{7. CloudWatch Log Group
(\texttt{main.tf:183-186})}{7. CloudWatch Log Group (main.tf:183-186)}}}

\begin{Shaded}
\begin{Highlighting}[]
\NormalTok{resource "aws\_cloudwatch\_log\_group" "logs" \{}
\NormalTok{  name              = "/ecs/newsrag{-}project"}
\NormalTok{  retention\_in\_days = 7}
\NormalTok{\}}
\end{Highlighting}
\end{Shaded}

\hypertarget{biux1ebfn-muxf4i-trux1b0ux1eddng-chung-main.tf189-218}{%
\subsection{\texorpdfstring{8. Biến môi trường chung
(\texttt{main.tf:189-218})}{8. Biến môi trường chung (main.tf:189-218)}}}

\emph{(Mã nguồn chi tiết đã được lược bỏ để báo cáo ngắn gọn. Vui lòng
xem trong source code dự án)}

\begin{quote}
\textbf{Lưu ý:} Trên AWS production, \texttt{QDRANT\_*} và
\texttt{KAFKA\_*} không dùng (thay bằng SQS + Aurora pgvector). Giữ lại
cho tương thích code local.
\end{quote}

\hypertarget{ecs-task-definitions-main.tf222-292}{%
\subsection{\texorpdfstring{9. ECS Task Definitions
(\texttt{main.tf:222-292})}{9. ECS Task Definitions (main.tf:222-292)}}}

Ba task definitions dùng chung Docker image:

\begingroup
\small
\setlength{\tabcolsep}{3pt}
\renewcommand{\arraystretch}{1.15}
\begin{longtable}[]{@{}
  >{\raggedright\arraybackslash}p{(\columnwidth - 6\tabcolsep) * \real{0.2143}}
  >{\raggedright\arraybackslash}p{(\columnwidth - 6\tabcolsep) * \real{0.1786}}
  >{\raggedright\arraybackslash}p{(\columnwidth - 6\tabcolsep) * \real{0.2857}}
  >{\raggedright\arraybackslash}p{(\columnwidth - 6\tabcolsep) * \real{0.3214}}@{}}
\toprule\noalign{}
\begin{minipage}[b]{\linewidth}\raggedright
Task
\end{minipage} & \begin{minipage}[b]{\linewidth}\raggedright
CPU
\end{minipage} & \begin{minipage}[b]{\linewidth}\raggedright
Memory
\end{minipage} & \begin{minipage}[b]{\linewidth}\raggedright
Command
\end{minipage} \\
\midrule\noalign{}
\endhead
\bottomrule\noalign{}
\endlastfoot
\textbf{crawler} & 256 (0.25 vCPU) & 512 MB &
\texttt{python\ main.py\ -\/-mode\ crawl} \\
\textbf{etl} & 512 (0.5 vCPU) & 1024 MB &
\texttt{python\ main.py\ -\/-mode\ etl} \\
\textbf{vectorize} & 512 & 1024 MB &
\texttt{python\ main.py\ -\/-mode\ vectorize} \\
\end{longtable}
\endgroup

Tất cả dùng: - \texttt{network\_mode\ =\ "awsvpc"} (bắt buộc cho
Fargate) - \texttt{execution\_role\_arn} = ECS execution role -
\texttt{task\_role\_arn} = ECS task role (Bedrock permissions) -
\texttt{awslogs} driver \(\rightarrow\) \texttt{/ecs/newsrag-project}

\hypertarget{eventbridge-scheduler-main.tf294-357}{%
\subsection{\texorpdfstring{10. EventBridge Scheduler
(\texttt{main.tf:294-357})}{10. EventBridge Scheduler (main.tf:294-357)}}}

Cron jobs hàng ngày (UTC):

\emph{(Mã nguồn chi tiết đã được lược bỏ để báo cáo ngắn gọn. Vui lòng
xem trong source code dự án)}

\textbf{Cron Expression Format:}
\texttt{cron(minutes\ hours\ day-of-month\ month\ day-of-week\ year)} -
\texttt{cron(0\ 1\ *\ *\ ?\ *)} = Hàng ngày lúc 01:00 UTC - \texttt{?} =
không có day-of-week cụ thể (dùng day-of-month thay thế)

\hypertarget{variables-main.tf5-38}{%
\subsection{\texorpdfstring{11. Variables
(\texttt{main.tf:5-38})}{11. Variables (main.tf:5-38)}}}

\begin{Shaded}
\begin{Highlighting}[]
\NormalTok{variable "db\_password" \{ type = string; sensitive = true \}}
\NormalTok{variable "db\_user" \{ type = string \}}
\NormalTok{variable "qdrant\_api\_key" \{ type = string; sensitive = true \}}
\NormalTok{variable "qdrant\_host" \{ type = string \}}
\NormalTok{variable "model\_1\_api\_key" \{ type = string; sensitive = true \}}
\NormalTok{variable "model\_2\_api\_key" \{ type = string; sensitive = true \}}
\NormalTok{variable "model\_3\_api\_key" \{ type = string; sensitive = true \}}
\end{Highlighting}
\end{Shaded}

Điền vào \texttt{terraform.tfvars} (không commit lên git):

\begin{Shaded}
\begin{Highlighting}[]
\NormalTok{db\_password     = "your\_secure\_password\_123!"}
\NormalTok{db\_user         = "postgres"}
\NormalTok{qdrant\_api\_key  = ""}
\NormalTok{qdrant\_host     = ""}
\NormalTok{model\_1\_api\_key = "gsk\_xxxxxxxxxxxxxxxxxxxx"  \# Groq}
\NormalTok{model\_2\_api\_key = "gsk\_xxxxxxxxxxxxxxxxxxxx"  \# Groq backup}
\NormalTok{model\_3\_api\_key = "AIza\_xxxxxxxxxxxxxxxxxxxx" \# Gemini}
\end{Highlighting}
\end{Shaded}

\hypertarget{outputs-main.tf360-370}{%
\subsection{\texorpdfstring{12. Outputs
(\texttt{main.tf:360-370})}{12. Outputs (main.tf:360-370)}}}

\begin{Shaded}
\begin{Highlighting}[]
\NormalTok{output "rds\_endpoint" \{ value = aws\_rds\_cluster.main.endpoint \}}
\NormalTok{output "ecr\_repository\_url" \{ value = aws\_ecr\_repository.api.repository\_url \}}
\NormalTok{output "ecs\_cluster\_name" \{ value = aws\_ecs\_cluster.cluster.name \}}
\end{Highlighting}
\end{Shaded}

Sau \texttt{terraform\ apply}, dùng outputs để: - Kết nối database:
\texttt{psql\ "postgresql://postgres:pass@\$(terraform\ output\ -raw\ rds\_endpoint):5432/newsrag"}
- Build \& push Docker:
\texttt{docker\ push\ \$(terraform\ output\ -raw\ ecr\_repository\_url):latest}
- Chạy task thủ công:
\texttt{aws\ ecs\ run-task\ -\/-cluster\ \$(terraform\ output\ -raw\ ecs\_cluster\_name)\ ...}

\hypertarget{triux1ec3n-khai}{%
\subsection{Triển khai}}

\emph{(Mã nguồn chi tiết đã được lược bỏ để báo cáo ngắn gọn. Vui lòng
xem trong source code dự án)}

\hypertarget{khux1edfi-tux1ea1o-database-schema}{%
\subsection{Khởi tạo Database
Schema}}

\begin{Shaded}
\begin{Highlighting}[]
\VariableTok{ENDPOINT}\OperatorTok{=}\VariableTok{$(}\ExtensionTok{terraform}\NormalTok{ output }\AttributeTok{{-}raw}\NormalTok{ rds\_endpoint}\VariableTok{)}
\ExtensionTok{psql} \StringTok{"postgresql://postgres:your\_password@}\VariableTok{$\{ENDPOINT\}}\StringTok{:5432/newsrag"} \AttributeTok{{-}f}\NormalTok{ database/warehouse.sql}
\end{Highlighting}
\end{Shaded}

File \texttt{database/warehouse.sql} tạo: - \textbf{Dimension Tables}:
\texttt{dim\_source}, \texttt{dim\_time}, \texttt{dim\_content},
\texttt{dim\_author} - \textbf{Fact Tables}: \texttt{fact\_articles},
\texttt{fact\_article\_authors}, \texttt{fact\_chunks} -
\textbf{pgvector extension + HNSW index} trên
\texttt{fact\_chunks.embedding} (vector(1024))

\hypertarget{xuxe1c-minh-triux1ec3n-khai}{%
\subsection{Xác minh triển khai}}

\emph{(Mã nguồn chi tiết đã được lược bỏ để báo cáo ngắn gọn. Vui lòng
xem trong source code dự án)}

\hypertarget{ux1b0ux1edbc-lux1b0ux1ee3ng-chi-phuxed-huxe0ng-thuxe1ng-ap-southeast-2}{%
\subsection{Ước lượng chi phí (Hàng tháng,
ap-southeast-2)}}

\begingroup
\small
\setlength{\tabcolsep}{3pt}
\renewcommand{\arraystretch}{1.15}
\begin{longtable}[]{@{}lll@{}}
\toprule\noalign{}
Resource & Configuration & Est. Cost \\
\midrule\noalign{}
\endhead
\bottomrule\noalign{}
\endlastfoot
Aurora Serverless v2 & 2 ACU avg (db.t4g.medium) &
\textasciitilde\$15-20 \\
ECS Fargate Crawler & 0.25 vCPU, 0.5 GB, 30 min/day &
\textasciitilde\$0.50 \\
ECS Fargate ETL & 0.5 vCPU, 1 GB, 30 min/day & \textasciitilde\$1.00 \\
ECS Fargate Vectorize & 0.5 vCPU, 1 GB, 30 min/day &
\textasciitilde\$1.00 \\
ECR & 1 repository, \textasciitilde2 GB storage &
\textasciitilde\$0.20 \\
CloudWatch Logs & 7-day retention, \textasciitilde1 GB ingest &
\textasciitilde\$1.00 \\
\textbf{Total (Infrastructure)} & &
\textbf{\textasciitilde\$19-24/month} \\
\end{longtable}
\endgroup

\begin{quote}
\textbf{Note:} Lambda + Bedrock + API Gateway costs (covered in later
sections) add \textasciitilde\$2-5/month. \textbf{Total
\textasciitilde\$21-29/month}.
\end{quote}

\hypertarget{dux1ecdn-dux1eb9p}{%
\subsection{Dọn dẹp}}

\begin{Shaded}
\begin{Highlighting}[]
\ExtensionTok{terraform}\NormalTok{ destroy}
\CommentTok{\# Nhập \textquotesingle{}yes\textquotesingle{} để xác nhận}
\end{Highlighting}
\end{Shaded}

\begin{quote}
! \textbf{CẢNH BÁO:} Lệnh này xóa Aurora cluster và \textbf{mất toàn bộ
dữ liệu vĩnh viễn}. Backup trước nếu cần.
\end{quote}

\begin{center}\rule{0.5\linewidth}{0.5pt}\end{center}

\textbf{Tiếp theo:} \href{5.4-Local-Dev/}{Thiết lập phát triển cục bộ}
